% GENERATED FILE. Edit canonical lessons, then run npm run generate:books.
% canonical-source-hash: fnv1a64:b35b2663d5344956
% canonical-lessons: ES-C773-cobertura, ES-C773-adaptador, ES-C773-aplicacion, ES-C773-enlace, ES-C773-recargar

\chapter{Conectado de viaje}
\label{ch:es-conectado-de-viaje}
\hlchaptermodality{773}

\begin{chapteropening}
\textbf{By the end of this chapter:} \emph{I can sort out my phone on a trip: the signal, a plug adapter, an app, a link, and topping up credit.}

Complete the last lesson of chapter 773: I can sort out my phone on a trip: the signal, a plug adapter, an app, a link, and topping up credit.
\end{chapteropening}

\section[la cobertura]{la cobertura — the signal, the coverage}

\label{lesson:ES-C773-cobertura}



\begin{quote}
Before the new one: say the Spanish for to slip, then the Spanish for the compass.
\end{quote}

\subsection*{You'll want to know: la cobertura}
\textbf{la cobertura} — \textquotedblleft{}the signal, the coverage\textquotedblright{}.

\begin{quote}
En el pueblo no hay cobertura.
\end{quote}

\begin{quote}
\emph{There is no signal in the village.}
\end{quote}

For an insurance policy it is the cover: \emph{la cobertura del seguro}.

\begin{cousinweb}
\emph{cobertura} comes from Late Latin \emph{coopertūra} 'a covering' (EtymDB), from \emph{coopertus}, the participle of \emph{cooperīre} 'to cover over', the verb behind \emph{cubrir}.
\end{cousinweb}

\begin{grammarlens}[title={What you've built}]
One more word for telling, explaining and describing at B1.
\end{grammarlens}

\subsection*{Your turn}
\begin{itemize}
\raggedright
  \item \emph{Say it:} \emph{la cobertura}
  \item \emph{Say it:} \emph{la cobertura}, once more
  \item \emph{Say it:} a sentence of your own with \emph{la cobertura}
  \item \emph{Recall:} say the Spanish for to slip, then the Spanish for the compass, then say \emph{la cobertura} again
\end{itemize}

\begin{tcolorbox}[breakable,colback=teal!4,colframe=teal!35!black,title={Before you move on}]
What does la cobertura mean? (\textquotedblleft{}The signal, the coverage\textquotedblright{}.) Say it once more.
\end{tcolorbox}
\section[el adaptador]{el adaptador — the adapter}

\label{lesson:ES-C773-adaptador}



\begin{quote}
Before the new one: say the Spanish for the compass, then the Spanish for the signal.
\end{quote}

\subsection*{You'll want to know: el adaptador}
\textbf{el adaptador} — \textquotedblleft{}the adapter\textquotedblright{}.

\begin{quote}
Necesito un adaptador para el enchufe del hotel.
\end{quote}

\begin{quote}
\emph{I need an adapter for the hotel socket.}
\end{quote}

\begin{cousinweb}
\emph{adaptador} is built on \emph{adaptar} with the ending \emph{-dor}. \emph{Adaptar} comes from Latin \emph{adaptāre} 'to fit to' (EtymDB), built on \emph{aptus} 'fitted'.
\end{cousinweb}

\begin{grammarlens}[title={What you've built}]
One more word for telling, explaining and describing at B1.
\end{grammarlens}

\subsection*{Your turn}
\begin{itemize}
\raggedright
  \item \emph{Say it:} \emph{el adaptador}
  \item \emph{Say it:} \emph{el adaptador}, once more
  \item \emph{Say it:} a sentence of your own with \emph{el adaptador}
  \item \emph{Recall:} say the Spanish for the compass, then the Spanish for the signal, then say \emph{el adaptador} again
\end{itemize}

\begin{tcolorbox}[breakable,colback=teal!4,colframe=teal!35!black,title={Before you move on}]
What does el adaptador mean? (\textquotedblleft{}The adapter\textquotedblright{}.) Say it once more.
\end{tcolorbox}
\section[la aplicación]{la aplicación — the app; the application}

\label{lesson:ES-C773-aplicacion}



\begin{quote}
Before the new one: say the Spanish for the signal, then the Spanish for the adapter.
\end{quote}

\subsection*{You'll want to know: la aplicación}
\textbf{la aplicación} — \textquotedblleft{}the app; the application\textquotedblright{}.

\begin{quote}
Descarga la aplicación del metro antes de salir.
\end{quote}

\begin{quote}
\emph{Download the metro app before you leave.}
\end{quote}

In speech it is often \emph{la app}.

\begin{cousinweb}
\emph{aplicación} is a learned borrowing from Latin \emph{applicātiō} 'a joining to' (EtymDB), from \emph{applicāre} 'to attach', built on \emph{plicāre} 'to fold'.
\end{cousinweb}

\begin{grammarlens}[title={What you've built}]
One more word for telling, explaining and describing at B1.
\end{grammarlens}

\subsection*{Your turn}
\begin{itemize}
\raggedright
  \item \emph{Say it:} \emph{la aplicación}
  \item \emph{Say it:} \emph{la aplicación}, once more
  \item \emph{Say it:} a sentence of your own with \emph{la aplicación}
  \item \emph{Recall:} say the Spanish for the signal, then the Spanish for the adapter, then say \emph{la aplicación} again
\end{itemize}

\begin{tcolorbox}[breakable,colback=teal!4,colframe=teal!35!black,title={Before you move on}]
What does la aplicación mean? (\textquotedblleft{}The app; the application\textquotedblright{}.) Say it once more.
\end{tcolorbox}
\section[el enlace]{el enlace — the link; the connection}

\label{lesson:ES-C773-enlace}



\begin{quote}
Before the new one: say the Spanish for the adapter, then the Spanish for the app.
\end{quote}

\subsection*{You'll want to know: el enlace}
\textbf{el enlace} — \textquotedblleft{}the link; the connection\textquotedblright{}.

\begin{quote}
Te mando el enlace del billete.
\end{quote}

\begin{quote}
\emph{I'll send you the link to the ticket.}
\end{quote}

At a station or airport it is a connection: \emph{perder el enlace}.

\begin{cousinweb}
\emph{enlace} is built on \emph{enlazar} 'to tie together' (EtymDB), which comes from Latin \emph{inlaqueāre}, built on \emph{laqueus} 'a noose' (Lewis \& Short).
\end{cousinweb}

\begin{grammarlens}[title={What you've built}]
One more word for telling, explaining and describing at B1.
\end{grammarlens}

\subsection*{Your turn}
\begin{itemize}
\raggedright
  \item \emph{Say it:} \emph{el enlace}
  \item \emph{Say it:} \emph{el enlace}, once more
  \item \emph{Say it:} a sentence of your own with \emph{el enlace}
  \item \emph{Recall:} say the Spanish for the adapter, then the Spanish for the app, then say \emph{el enlace} again
\end{itemize}

\begin{tcolorbox}[breakable,colback=teal!4,colframe=teal!35!black,title={Before you move on}]
What does el enlace mean? (\textquotedblleft{}The link; the connection\textquotedblright{}.) Say it once more.
\end{tcolorbox}
\section[recargar]{recargar — to top up; to recharge}

\label{lesson:ES-C773-recargar}



\begin{quote}
Before the new one: say the Spanish for the app, then the Spanish for the link.
\end{quote}

\subsection*{You'll want to know: recargar}
\textbf{recargar} — \textquotedblleft{}to top up; to recharge\textquotedblright{}.

\begin{quote}
Tengo que recargar el móvil antes del viaje.
\end{quote}

\begin{quote}
\emph{I have to top up my phone before the trip.}
\end{quote}

It takes \emph{gu} before \emph{e}: \emph{recargué}. It also means to recharge a battery.

\begin{cousinweb}
\emph{recargar} is built on \emph{cargar} 'to load, to charge' with \emph{re-} 'again'. \emph{Cargar} comes from Late Latin \emph{carricāre} 'to load a cart', from \emph{carrus} 'a cart' (Lewis \& Short).
\end{cousinweb}

\begin{grammarlens}[title={What you've built}]
One more word for telling, explaining and describing at B1.
\end{grammarlens}

\subsection*{Your turn}
\begin{itemize}
\raggedright
  \item \emph{Say it:} \emph{recargar}
  \item \emph{Say it:} \emph{recargar}, once more
  \item \emph{Say it:} a sentence of your own with \emph{recargar}
  \item \emph{Recall:} say the Spanish for the app, then the Spanish for the link, then say \emph{recargar} again
\end{itemize}

\begin{tcolorbox}[breakable,colback=teal!4,colframe=teal!35!black,title={Before you move on}]
What does recargar mean? (\textquotedblleft{}To top up; to recharge\textquotedblright{}.) Say it once more.
\end{tcolorbox}
